% MAS1004 Assignment 2: long report (written with my agent).
% Every number below comes from experiments.csv, results/*/result.json or the
% saved outputs in outputs/, all printed by code in this repository.

\documentclass[11pt]{article}
\usepackage{mas1004}

\reporttitle{Assignment 2: Long report}
\studentname{이다윗}
\studentid{20231800}

\begin{document}
\makeheader

리더보드 닉네임: \texttt{rokafsogang}\quad 레포: \url{https://github.com/rokafsogang/mas1004-assignment2}\\
마지막 릴리즈: Track A \texttt{track-a-3} (balanced + 사진 115장 제외), Track B \texttt{track-b-2} (MobileNetV3-Large, 384px).

\textbf{학습 환경.} Track A는 MacBook Air(M5)의 GPU(\texttt{mps})에서, Track B는 Colab GPU(\texttt{cuda})에서 학습했다.
torch 2.6은 \texttt{mps}에서 \texttt{official.py}가 요청하는 \texttt{channels\_last} 메모리 형식 때문에 첫 backward에서
``view size is not compatible with input tensor's size and stride'' 오류로 멈췄다. \texttt{official.py}는 고치지 않고,
\texttt{.to(memory\_format=torch.channels\_last)}만 아무 일도 하지 않게 만든 래퍼 \texttt{scripts/official\_mps.py}로 실행했다.
recipe, 사진 선택, 예산(249,984장), seed, 검증 방식은 \texttt{official.py} 그대로이고, 텐서의 메모리 배치만 다르다.
채점기는 \texttt{cuda}에서 float16으로 학습하므로 내 검증 점수와 채점기 점수는 조금 다를 수 있다.

% ---------------------------------------------------------------------------
\section*{Track A}

\begin{center}\footnotesize
\begin{tabular}{lp{0.58\linewidth}r}
\toprule
tag & 이전 실행에서 바뀐 것 & validation \\
\midrule
baseline (seed 0) & starter recipe 그대로, 학습 사진 44,257장 전부 & 0.4961 \\
baseline (seed 1) & seed만 0 $\to$ 1 (우연에 의한 흔들림 측정) & 0.4956 \\
balanced & \texttt{sample\_weights}: 각 사진에 1/(그 도시 사진 수), 도시를 같은 빈도로 뽑음 & 0.5076 \\
balanced\_drop & 사진을 보고 고른 115장 제외 (44,257 $\to$ 44,142장), recipe는 balanced 그대로 & 0.5099 \\
no\_flip & \texttt{T.RandomHorizontalFlip()} 제거, 사진 선택은 balanced\_drop 그대로 & 0.5001 \\
\bottomrule
\end{tabular}
\end{center}

\begin{center}\footnotesize
\begin{tabular}{lrrrrrr}
\toprule
도시별 정답률 & seoul & busan & tokyo & osaka & taipei & 평균 \\
\midrule
baseline & 0.595 & 0.386 & 0.624 & 0.344 & 0.531 & 0.496 \\
balanced & 0.496 & 0.471 & 0.520 & 0.435 & 0.616 & 0.508 \\
balanced\_drop & 0.536 & 0.459 & 0.517 & 0.423 & 0.614 & 0.510 \\
no\_flip & 0.550 & 0.430 & 0.511 & 0.419 & 0.590 & 0.500 \\
\midrule
학습 사진 수 & 11,274 & 5,191 & 13,347 & 7,200 & 7,245 & \\
\bottomrule
\end{tabular}
\end{center}

\textbf{seed에 따른 흔들림.} 같은 starter recipe를 seed 0과 1로 학습하니 0.4961과 0.4956으로 평균은 0.0005 차이였다. 그러나 도시별로는 오사카 0.344 $\to$ 0.367(+0.024), 도쿄 0.624 $\to$ 0.607($-$0.017)처럼 더 크게 움직였다. 결과를 어떻게 보는지(내 말): ``같은 설정인데도 값이 살짝 다른걸 보면 기계라고 그 전껄 100\%같게 판독하지 못한다는 걸 알았다. 사람처럼 얘네도 시행마다 살짝씩 다르단걸''.
README에 따르면 같은 recipe도 다시 학습하면 2--3점 움직일 수 있다. 아래의 각 변경은 seed 0으로 한 번씩만 학습했으므로,
이 흔들림보다 작은 차이는 우연일 수 있다.

\textbf{balanced.} 학습 사진 수가 도쿄 13,347장, 부산 5,191장으로 2.6배 차이가 난다. 왜 골랐는지(내 말):
``어차피 최대 양의 사진 (도쿄)꺼를 해야하는데 그거만 13,347할게 아니라 그 시간에 부산도 여러번 보면 정확도가 올라갈거라 생각해서 밸런스 선택''.
평균은 0.496에서 0.508로 올랐다. 부산(+0.085), 오사카(+0.091), 타이베이(+0.085)는 올랐고 서울($-$0.099), 도쿄($-$0.104)는 내려갔다.
검증 사진 12,077장에 대한 모델의 답을 세어 보면, baseline은 ``부산''이라고 938번 답했고(실제 부산 사진 1,390장) ``도쿄''라고 4,320번 답했다(실제 3,881장).
balanced는 ``부산'' 1,362번, ``도쿄'' 3,318번으로, 사진이 많은 도시 쪽으로 쏠리던 답이 고르게 바뀌었다
(\texttt{results/*/result.json}의 confusion matrix). 결과를 어떻게 보는지(내 말):
``서울 도쿄를 제외한 나머지 도시는 베이스라인보다 높은 정확도를 보여줬는데 그 이유는 아무래도 사진 양이 적었기에 여러번 학습하여 정확도가 올라갔다고 보여짐''.

\textbf{no\_flip.} 좌우반전을 하면 간판 글자가 거울 글씨가 되고 차가 다니는 쪽이 바뀐다. 왜 골랐는지(내 말):
``좌우반전을 하면 너 말대로 일본과 같은 반대인 나라에 대한 오류가 생길거라 생각이 들어''.
다른 후보(더 강한 확대 자르기, label smoothing 끄기)를 고르지 않은 이유(내 말):
``확대해봤자 픽셀이 더 세분화되는거 아니니 도움 안될거같고 100\%라는건 세상에 없어''.
결과: 0.5099 $\to$ 0.5001($-$0.0098)로 내려갔다. 좌측통행인 도쿄($-$0.007)와 오사카($-$0.004)도 오르지 않았고, 부산($-$0.029)과 타이베이($-$0.024)가 가장 많이 내려갔으며 서울만 +0.014 올랐다. 결과를 어떻게 보는지(내 말): ``예상과 달리 떨어졌고, 이유는 모르겠다''. 그래서 릴리즈하지 않았고 \texttt{track\_a/recipe.py}는 track-a-3의 버전으로 되돌렸다.

\textbf{최종 recipe와 리더보드.} 마지막 Track A 릴리즈는 \texttt{track-a-3}(balanced + 115장 제외, validation 0.5099)이다.
리더보드 점수는 이 보고서를 쓴 시점(11월 9일 밤)에 아직 공개되지 않았다. 매일 자정 채점 결과가 다음 날 아침에 올라오고,
나는 11월 9일에 처음 릴리즈했기 때문이다. 그래서 검증 점수와 리더보드 점수를 나란히 비교하지 못했다.

% ---------------------------------------------------------------------------
\section*{사진을 보고 내린 결정}

\picturefile{figures/suspects_baseline_seed0_busan_00.png}

\textbf{무엇을 봤나.} baseline 모델로 \texttt{curate.py suspects results/baseline\_seed0/model.onnx --every 4}를 돌려,
각 도시 학습 사진 중 모델이 그 도시일 확률을 가장 낮게 준 48장씩(위 그림은 부산)을 봤다.
처음 도시별 사진(\texttt{curate.py look})을 봤을 때 도시를 구분하는 단서로 적은 것(내 말):
``서울 잠실 주경기장을 통해 알수 있음, 부산 - 해운대 모래사장을 통해 알 수 있음, 도쿄 - JR전철을 통해(초록색), 오사카 - 도톤보리 수로를 통해, 타이페이 - 한자와 타이페이101을 통해 알 수 있음''.

\textbf{규칙 (빼기 전에 정함, 내 말).} ``사람사진과 음식 사진 그 두개는 어디 나라인지정도까진 추측이 가능해도 도시를 알기는 어려움 그리고 의미없는 자동차 이런거 다 없앰''.

\textbf{무엇을 바꿨나.} 이 규칙에 맞는 사진을 시트 번호로 직접 골라 115장을 뺐다(서울 25, 부산 29, 도쿄 26, 오사카 25, 타이베이 10).
\texttt{curate.py count}의 출력:
\outputfile{outputs/curate_count.txt}

\textbf{점수.} balanced 0.5076 $\to$ balanced\_drop 0.5099 (+0.0023). 서울은 0.496 $\to$ 0.536으로 올랐고 나머지 도시는 거의 그대로이거나 조금 내려갔다.
115장은 학습 사진의 0.26\%이고, 차이는 seed 흔들림보다 작을 가능성이 크다. 결과를 어떻게 보는지(내 말):
``아무래도 그 많은 사진 중에 115장 밖에 안뺐는데 당연히 정확도가 거의 그대로겠지 정말 다량의 사진을 검수 수거 하지 않는한 크게 안바뀔거야''.

\textbf{에이전트에게 준 요청.} 사진을 고른 것과 규칙은 내가 정했고, 에이전트는 그 번호를 그대로 \texttt{curate.py drop}으로 뺐다.
\outputfile{outputs/requests.txt}

% ---------------------------------------------------------------------------
\section*{Track B}

\begin{center}\footnotesize
\begin{tabular}{lp{0.58\linewidth}r}
\toprule
tag & 바뀐 것 & validation \\
\midrule
b1\_mnv3l\_320 & 첫 모델: MobileNetV3-Large(\texttt{IMAGENET1K\_V2}), 320px, 6 epoch, 도시 균등(\texttt{--balance}), starter 사진만 & 0.5631 \\
b2\_mnv3l\_384 & 이미지 크기만 320 $\to$ 384 & 0.5699 \\
\bottomrule
\end{tabular}
\end{center}

\textbf{b1.} \texttt{train\_b.py} 기본값(ResNet18, 224px, 4 epoch)과 비교하면 모델, 이미지 크기, epoch, 균등 뽑기를 한꺼번에 바꿨기 때문에,
이 중 무엇이 도움이 됐는지는 b1만으로는 알 수 없다. 4.20M 파라미터, 0.44 GMACs.
Track A의 balanced(0.508)보다 높게 나온 이유를 어떻게 보는지(내 말):
``사진을 크게 보면 픽셀이 조각되면서 컴퓨터가 학습을 조금더 세밀하게할 수 있어서 더 높은 값을 추출했다''.

\textbf{b2.} 위의 생각을 확인하려고 이미지 크기 하나만 384로 키웠다(0.63 GMACs). 0.5631 $\to$ 0.5699(+0.0068)로 올랐다.
서울(+0.032), 도쿄(+0.023)는 올랐고 부산($-$0.017), 오사카($-$0.017)는 내려갔다. 결과를 어떻게 보는지(내 말): ``픽셀을 더 쪼개야한다고 봐''.
seed마다 다시 학습해 보지 않았으므로 +0.007이 크기 때문인지 우연인지는 확실하지 않다.

\textbf{추가 사진.} 모으지 않았다. 두 모델 모두 starter 학습 사진만 썼다.

\textbf{검증 점수와 리더보드 점수.} b1은 11월 9일 \texttt{track-b-1}, b2는 같은 날 23:37에 \texttt{track-b-2}로 릴리즈했다.
리더보드 점수는 보고서 작성 시점에 아직 공개되지 않았다. 두 점수가 다를 수 있는 이유: 검증 사진은 Wikimedia Commons, OSV-5M, YouTube가 섞인
starter와 같은 출처의 사진이지만, 리더보드 사진은 starter에 없는 채널의 YouTube 거리 영상 프레임이고, 최종 사진은 휴대폰으로 찍은 거리 사진이다.
학습 사진의 절반 이상이 Wikimedia Commons 사진이라, 거리 영상 프레임에서는 점수가 더 낮게 나올 수 있다.

\subsection*{채점기와 같은 방식으로 본 제출물 (b2)}
\outputfile{outputs/check_submission.txt}

% ---------------------------------------------------------------------------
\section*{무엇이 차이를 만들었나}

\begin{itemize}
  \item Track A에서 가장 큰 변화는 balanced(+0.0115)였다. 사진을 빼기(+0.0023)는 거의 변화가 없었다. no\_flip은 $-$0.0098로 오히려 내려갔다.
  \item Track B에서 이미지 크기 320 $\to$ 384는 +0.0068이었다.
  \item 두 트랙 사이의 차이(Track A 최고 0.510, Track B 최고 0.570)가 가장 컸지만, Track B는 모델, 이미지 크기, epoch을 함께 바꿨으므로 어느 것 때문인지는 따로 재지 않았다.
  \item 각 변경은 seed 0으로 한 번씩만 학습했다. baseline의 두 seed는 평균이 0.0005밖에 차이 나지 않았지만 도시별로는 0.024까지 움직였고, 흔들림을 두 번만 쟀으므로 balanced(+0.0115)처럼 큰 차이와 달리 115장 제외(+0.0023) 같은 작은 차이는 우연과 확실히 구분할 수 없다.
\end{itemize}

\textbf{일주일이 더 있다면(내 말).} ``픽셀을 훨씬 더 많이 줄여서 이걸 반복해서 실험해보고 싶음 내 가설 증명하고 싶음''.

% ---------------------------------------------------------------------------
\section*{All runs}
\outputfile{outputs/experiments.csv}

\section*{Lab notebook}
\outputfile{outputs/LOG.md}

\end{document}
